\section{Certificates and algorithmic boundaries}
\label{sec:algorithms}

All algorithms in this section take exact rational input. The input length
includes the binary encodings of its numerators and denominators. A
positive semidefinite matrix of small rank need not have a rational Gram
factor with the same row dimension as its rank, so we work with a
rational quotient and its rational inner product.

\subsection{Rational linear algebra and short certificates}

\begin{lemma}\label{lem:rational-negative-direction}
For a rational symmetric matrix $A$, polynomial-time exact linear algebra
either certifies $A\succeq0$ and finds a positive definite principal
submatrix $A_{KK}$ of order $\rank A$, or returns an integer vector $y$ of
polynomial encoding length with $y\trans Ay<0$.
\end{lemma}

\begin{proof}
Apply symmetric elimination to the current Schur complement $C$. A
negative diagonal entry gives a negative direction $e_i$. If $C_{ii}=0$
and $C_{ij}\ne0$, take $z_j=1$ and
$z_i=-(C_{jj}+1)/(2C_{ij})$, with all other entries zero; then
$z\trans Cz=-1$. Otherwise pivot on a positive diagonal entry. For a
partition with that entry first, a direction $z_R$ in the new Schur
complement lifts by setting $z_i=-C_{iR}z_R/C_{ii}$, without changing its
quadratic value. If the remaining matrix is zero, the matrix is positive
semidefinite, and the positive pivots identify the required set $K$.

After pivots indexed by $K$, the Schur-complement entries are ratios of
minors of the original matrix. Clearing input denominators and applying
the determinant bound shows that these ratios have polynomial encoding
length. A negative direction can be lifted at once by
$z_K=-A_{KK}^{-1}A_{KR}z_R$, where $z_R$ has at most two nonzero entries.
Cramer's rule gives polynomial encoding length for the lifted direction.
Multiplying by its common denominator gives the asserted integer vector.
Fraction-free elimination, or rational elimination with reduced
fractions, performs these operations in polynomial time.
\end{proof}

\begin{lemma}\label{lem:rational-quotient}
Let $A\succeq0$ be rational of rank $r>0$. In polynomial time one can find
$K$, $P=A_{K,\cdot}$ and $G=A_{KK}^{-1}\succ0$ such that
\[
 A=P\trans GP,
\]
and a rational basis $R_\Lambda$ of the rank-$r$ lattice
$\Lambda=P\Z^m$, together with an integer matrix $T$ of polynomial
encoding length satisfying $R_\Lambda=PT$. If a rational vector $y$ of polynomial
encoding length belongs to $\Lambda$, an integer solution of $Pz=y$ of
polynomial encoding length can be found in polynomial time.
\end{lemma}

\begin{proof}
Choose $K$ by \cref{lem:rational-negative-direction}. The Schur complement
of $A_{KK}$ is zero because $A$ has rank $r$, which gives $A=P\trans GP$.
Let $D$ be a common denominator of $P$. Its encoding length is at most the
sum of the encoding lengths of the denominators of $P$. The integer
matrix $DP$ has full row rank. Its column Hermite normal form gives a
basis $H$ of $(DP)\Z^m$ and solves integral systems with this coefficient
matrix in polynomial time, with polynomial encoding lengths
\citep{KannanBachem1979}. Take $R_\Lambda=H/D$, and solve
$(DP)t_j=H_{\cdot j}$ for each basis column to obtain $T$.

Thus $\Lambda=R_\Lambda\Z^r$ is a lattice in $\R^r$. For any $y\in\Lambda$,
solving $(DP)z=Dy$ gives the final assertion. Here $Dy$ is integral; all
input and output lengths in this system are polynomial. This argument
also explains why linear dependence among the columns of a real Gram
factor causes no difficulty.
\end{proof}

\begin{lemma}\label{lem:short-hypermetric-witness}
For any rational symmetric $M$, if $g_M(z)<0$ for some $z\in\Z^m$, there
is such a vector of encoding length polynomial in the encoding length of
$M$. Consequently, hypermetric violation has polynomial certificates.
\end{lemma}

\begin{proof}
Put $\delta=\diag M$. If $M$ is not positive semidefinite, take the integer
negative direction from \cref{lem:rational-negative-direction}, and choose
its sign so that $\delta\trans z\ge0$. Then
$g_M(z)=z\trans Mz-\delta\trans z<0$.

Suppose next that $M\succeq0$ but $\delta\notin\range M$. Since
$\range M=(\ker M)^\perp$, a rational kernel basis contains a vector $k$
with $\delta\trans k\ne0$. Exact elimination supplies such a vector with
polynomial encoding length. Clear denominators and choose its sign so
that $\delta\trans k>0$. Then $g_M(k)<0$.

It remains to consider $M\succeq0$ and $\delta\in\range M$. If $M=0$,
then $\delta=0$ and there is no violator. Otherwise use
\cref{lem:rational-quotient}, writing $D_0=M_{KK}$, $P=M_{K,\cdot}$ and
$G=D_0^{-1}$. The range condition gives
$\delta=P\trans G\delta_K$. With $y=Pz$ and $t=\delta_K/2$,
\begin{equation}\label{eq:hypermetric-cvp}
 g_M(z)=(y-t)\trans G(y-t)-t\trans Gt.
\end{equation}
A violator therefore satisfies $\norm{y}_G<2\norm{t}_G$, where
$\norm{a}_G^2=a\trans Ga$. By Cauchy--Schwarz,
\[
 |y_i|\le \sqrt{(D_0)_{ii}}\,\norm{y}_G
 <2\sqrt{(D_0)_{ii}}\,\norm{t}_G.
\]
The square of the bound is rational of polynomial encoding length, so
its magnitude is at most $2^{\poly(L)}$, where $L$ is the input length.
Moreover, every coordinate of $y=Pz$ has denominator dividing a common
denominator of $P$. Hence $y$ itself has polynomial encoding length.
\Cref{lem:rational-quotient} gives a polynomial-size integer $z'$ with
$Pz'=y$. Equation~\eqref{eq:hypermetric-cvp} gives $g_M(z')=g_M(z)<0$.
The corresponding hypermetric vector is $(1-\one\trans z',z')$, whose
encoding length is also polynomial.
\end{proof}

For general integer moments, let
$Y=\left[\begin{smallmatrix}1&x\trans\\x&B\end{smallmatrix}\right]$
be rational and symmetric. No Boolean diagonal identities are assumed
unless stated. We use the parity identity \eqref{eq:parity-split},
with $u=e_0+2v$.

\begin{lemma}\label{lem:short-split-witness}
If a rational symmetric $Y$ with $Y_{00}=1$ violates an integer split,
then it has a violating integer split vector of polynomial encoding
length. Thus unrestricted integer split violation belongs to $\NP$.
\end{lemma}

\begin{proof}
If $Y$ is not positive semidefinite, take an integer negative direction
$v$ from \cref{lem:rational-negative-direction} and choose its sign so
that $v\trans Ye_0\le0$. Then $q_Y(v)<0$.

Otherwise apply \cref{lem:rational-quotient} to $Y$. A violated split gives
$u=e_0+2v$ with $u\trans Yu<1$. Set $y=Pu$. Since $Y=P\trans GP$,
$y\trans Gy<1$, and Cauchy--Schwarz gives
\[
 |y_i|<\sqrt{(Y_{KK})_{ii}}.
\]
The common denominator of $P$ bounds the denominators of $y$, so $y$ has
polynomial encoding length. The integral system
$2Pv'=y-Pe_0$ is feasible because $v$ solves it. After clearing
denominators, Hermite normal form produces a polynomial-size integer
solution $v'$. Then $P(e_0+2v')=y$, and
\eqref{eq:parity-split} gives $q_Y(v')<0$.
\end{proof}

The positive definite binary points in \cref{cor:hypermetric-hardness}
are also general integer moment inputs. Their reduction, together with
\cref{lem:short-split-witness}, proves strong $\NP$-completeness of
unrestricted integer split violation.

\subsection{Fixed rank, support and gonality}

\begin{theorem}\label{thm:rank-separation}
Hypermetric separation for rational $M$ of rank $r$, and integer split
separation for rational $Y$ of rank $r$ with $Y_{00}=1$, can be solved in
time $r^{O(r)}\poly(L)$, where $L$ is the input encoding length. The
algorithm either returns a violated inequality or certifies that none
exists; rank-zero inputs are decided directly.
\end{theorem}

\begin{proof}
For hypermetric separation the first two cases in the proof of
\cref{lem:short-hypermetric-witness} are polynomial-time constructive.
In the remaining case, \eqref{eq:hypermetric-cvp} asks whether the
distance from $t$ to $\Lambda=P\Z^m$, in the inner product $G$, is less
than $\norm{t}_G$. Construct $R_\Lambda=PT$ by
\cref{lem:rational-quotient}. The problem is exact closest vector in
dimension $r$, with rational Gram matrix $F=R_\Lambda\trans G R_\Lambda$ and rational
target coordinates $\xi=R_\Lambda^{-1}t$. To give a rational Euclidean instance,
compute a rational factorization $F=L_0\diag(d_i)L_0\trans$ with
$d_i=p_i/q_i>0$. The integer $p_iq_i$ is a sum of polynomially many
integer squares: in its binary expansion an even-position bit is one
square, and an odd-position bit is two equal squares. Dividing the
square roots by $q_i$ expresses $d_i$ as a sum of rational squares.
If $d_i=\sum_j\alpha_{ij}^2$, take the rows of $C$ to be
$\alpha_{ij}e_i\trans L_0\trans$.  This rectangular rational matrix
has polynomial encoding length and $C\trans C=F$. Thus our problem is
Euclidean closest vector to $C\xi$ in the rank-$r$ lattice $C\Z^r$.
The exact fixed-dimensional closest-vector algorithm takes
$r^{O(r)}\poly(L)$ time \citep[Theorem~4.5]{Kannan1987}. Its lattice
steps use inner products, projections and squared norms of the $r$
generators. These can equally be computed from their rational Gram
matrix; a larger ambient dimension affects only the polynomial cost of
forming inner products. A returned coefficient
vector $a\in\Z^r$ lifts to the hypermetric witness $z=Ta$.

For split separation, non-positive-semidefinite inputs are handled by
\cref{lem:rational-negative-direction}. Otherwise put $y_0=Pe_0$.
Then $y_0\trans Gy_0=Y_{00}=1$, and
\[
 q_Y(v)<0
 \quad\Longleftrightarrow\quad
 \norm{Pv+y_0/2}_G^2<\frac14.
\]
This is exact closest vector to $-y_0/2$ in the same rank-$r$ lattice.
A returned coefficient vector $a$ yields $v=Ta$. All quotient data have
polynomial encoding length. A nearest lattice point has norm at most
twice the target norm, so its rational coordinates, and then its integer
basis coordinates, have polynomial encoding length as well. Exact
comparison treats equality at the threshold as satisfaction, as required
by strict separation.
\end{proof}

\begin{corollary}\label{cor:bounded-support-separation}
For every fixed $k$, hypermetric and rounded psd separation restricted to
coefficient vectors supported on at most $k$ vertices is polynomial. Integer split
separation restricted to $|\supp w|\le k$ is also polynomial, with no
bound on the integer coefficients. For every fixed $k$, separation over
hypermetric or rounded psd inequalities of gonality at most $k$ is
polynomial.
\end{corollary}

\begin{proof}
For hypermetrics, enumerate all nonempty vertex sets of size at most $k$ and apply
\cref{thm:rank-separation} to each induced distance vector, choosing any
vertex of that set as root. Its correlation matrix has order at most
$k-1$, so the running time per set is polynomial. Every supported
hypermetric inequality belongs to one such induced instance. This avoids
any need for a bound on facet coefficients.
For rounded psd inequalities, apply the split correspondence to the
Boolean moment matrix of each induced distance vector and use the split
part of \cref{thm:rank-separation}; the matrix has order at most $k$.

For splits, enumerate sets $I$ of at most $k$ variable indices and apply
\cref{thm:rank-separation} to $Y_{\{0\}\cup I,\{0\}\cup I}$. Its rank is at
most $k+1$. There are $O(n^k)$ such sets.

Finally, a coefficient vector with $\sum_i|b_i|\le k$ is a sum of at most
$k$ signed unit vectors. Enumerating these descriptions gives
$O((2|V|+1)^k)$ candidates; filter by the defining parity or sum condition
and evaluate each inequality exactly. These support and gonality bounds
give polynomial algorithms for fixed parameters. Their displayed
dependence on the number of vertices is not a fixed-parameter tractability
claim for support or gonality.
\end{proof}

\subsection{A fixed normalized threshold}

Put $S=B-xx\trans$ and define the normalized violation of a nonzero
direction $w$ and split vector $v=(v_0,w)$ as
$-q_Y(v)/\norm{w}^2$. The denominator is the squared Euclidean norm of
the variable direction.

\begin{lemma}\label{lem:best-split-offset}
For $t=w\trans x$, write $\{t\}=t-\lfloor t\rfloor$ for its fractional
part and put $\phi(t)=\{t\}(1-\{t\})$. Then
\[
 \max_{v_0\in\Z}(-q_Y(v_0,w))
 =\phi(t)-w\trans Sw.
\]
If $Y\succeq0$, this quantity is at most $1/4$.
\end{lemma}

\begin{proof}
Expansion gives
$q_Y(v_0,w)=(v_0+t)(v_0+t+1)+w\trans Sw$.
Writing $f=\{t\}$, the minimum of the first term is $-f(1-f)$,
attained at $v_0=-\lfloor t\rfloor-1$. When $f=0$, the offset
$v_0=-\lfloor t\rfloor$ also attains it. If $Y\succeq0$, its Schur
complement $S$ is positive semidefinite.
\end{proof}

\begin{theorem}\label{thm:threshold-separation}
Let $Y\succeq0$ be rational with $Y_{00}=1$, and let $\rho>0$ be rational.
An exact algorithm finds a split with normalized violation greater than
$\rho$, or certifies that none exists, in time
\[
 f(B_\rho)n^{B_\rho}\poly(L+\operatorname{size}(\rho)),
 \qquad B_\rho=\max\left\{0,
             \left\lceil\frac1{4\rho}\right\rceil-1\right\}.
\]
In particular the problem is polynomial for each fixed positive
$\rho$. The bound is XP in the inverse threshold.
\end{theorem}

\begin{proof}
By \cref{lem:best-split-offset}, any qualifying direction satisfies
\[
 0<\norm{w}^2<\frac1{4\rho},
 \qquad\text{hence}\qquad \norm{w}^2\le B_\rho.
\]
If $B_\rho=0$, there is no qualifying direction. Otherwise enumerate all
integer directions with this bound. Such a direction has at most
$B_\rho$ nonzero entries, each of magnitude at most
$\lfloor\sqrt{B_\rho}\rfloor$. The number of candidates is at most
\[
 \sum_{j=1}^{\min(n,B_\rho)}
 \binom nj(2\lfloor\sqrt{B_\rho}\rfloor)^j
 \le (B_\rho+1)(2n\sqrt{B_\rho})^{B_\rho}.
\]
For each candidate compute the best offset by
\cref{lem:best-split-offset}, and compare
$\phi(w\trans x)-w\trans Sw$ with $\rho\norm{w}^2$ in exact rational
arithmetic. The candidate coefficients have $O(\log(B_\rho+1))$ bits;
the dot products, floor, fractional part and comparison have polynomial
bit complexity in their input lengths. The best offset also has
polynomial encoding length. Thus this finite search is exact, including
at equality with the threshold.
\end{proof}

The same argument bounds any adaptive search after a positive normalized
violation has been found. The theorem requires positive semidefinite
input. A polynomial algorithm that finds an arbitrary split at a
non-positive-semidefinite input does not by itself solve a normalized
threshold problem there. Nor does the theorem establish an approximation
hardness result for normalized violation.

\begin{proposition}\label{prop:raw-quarter}
For rational $Y\succeq0$ with $Y_{00}=1$, the maximum raw split violation
$\mu(Y)=\max_{v\in\Z^{n+1}}(-q_Y(v))$ is attained and lies in $[0,1/4]$.
It equals $1/4$ if and only if the integer kernel of $Y$ meets
$e_0+2\Z^{n+1}$. This equality condition is decidable in polynomial time.
\end{proposition}

\begin{proof}
The parity identity gives $q_Y(v)\ge-1/4$, and $v=0$ gives zero.
For a common denominator $D$ of $Y$, the nonnegative values
$u\trans Yu$, $u\in e_0+2\Z^{n+1}$, belong to $D^{-1}\Z_{\ge0}$.
This nonempty discrete set has a minimum, proving attainment.
The maximum is $1/4$ exactly when $u\trans Yu=0$ for some such $u$.
Positive semidefiniteness makes this equivalent to $Yu=0$.
After clearing denominators, Hermite normal form decides feasibility of
the integral system $2Yv=-Ye_0$ in polynomial time and supplies a split
attaining $1/4$ whenever the system is feasible.
\end{proof}

\begin{corollary}\label{cor:raw-approximation}
Unless $\mathrm P=\NP$, the maximum raw split violation at rational
positive definite Boolean moment points has no polynomial-time
multiplicative approximation with any finite factor.  This means a
value $\widehat\mu$ satisfying
$\mu(Y)/\alpha\leq\widehat\mu\leq\alpha\mu(Y)$ for finite
$\alpha\geq1$, including when $\mu(Y)=0$.
There is also no polynomial-time algorithm guaranteed to approximate
this value with absolute error smaller than $1/(4N)$ on the points
$Y_\varepsilon$ of \cref{thm:relaxation-hardness}.
\end{corollary}

\begin{proof}
The complete split classification in \cref{thm:relaxation-hardness}
gives $\mu(Y_\varepsilon)=1/(2N)$ when the X3C instance has a cover,
and $\mu(Y_\varepsilon)=0$ otherwise.  A finite-factor guarantee
distinguishes these cases by whether $\widehat\mu$ is positive.
An absolute error strictly smaller than $1/(4N)$ distinguishes them
by comparison with $1/(4N)$.  The input and this requested accuracy
have polynomial encoding length.  Both algorithms would decide X3C.
The reduction does not give hardness at a fixed additive tolerance.
\end{proof}

An integer direction $w\ne0$ is primitive if the greatest common
divisor of the absolute values of its coefficients is one.

\begin{lemma}\label{lem:primitive-split-domination}
Let $w=kw'$ with $0\ne w'\in\Z^n$ and integer $k\ge2$, and let
$v=(-s-1,w)$ with $s\in\Z$. Put
\begin{align*}
 t&=\left\lfloor\frac{2s+1-k}{2k}\right\rfloor,
 &b&=\frac{(2s+1)k-(2t+1)k^2}{2},\\
 a&=k^2-b,
 &c&=(k(t+1)-s)(k(t+1)-s-1).
\end{align*}
Then $a,b,c\ge0$, $a+b=k^2$, and for every symmetric $Y$,
\[
 q_Y(-s-1,kw')=
 a q_Y(-t-1,w')+b q_Y(-t-2,w')+cY_{00}.
\]
If $Y_{00}\ge0$, one of the two component splits has normalized violation
at least that of $v$. In particular primitive directions suffice for
violation detection and normalized optimization.
\end{lemma}

\begin{proof}
Write $z=(w')\trans Y_{\{1,\ldots,n\},0}$ and
$Z=(w')\trans Y_{\{1,\ldots,n\},\{1,\ldots,n\}}w'$. The left-hand side is
$k^2Z-(2s+1)kz+s(s+1)Y_{00}$. Matching the coefficients of $Z$ and $z$
with the two component slacks gives the stated $a$ and $b$. To obtain
the remaining constant, apply the same identity to
$g(z)=(kz-s)(kz-s-1)$ and $f_j(z)=(z-j)(z-j-1)$.
The quadratic and linear coefficients already match, and both
$f_t(t+1)$ and $f_{t+1}(t+1)$ vanish. Thus the constant is
$c=g(t+1)$. It is nonnegative because it is a product of two consecutive
integers. The definition of $t$ gives
$(2t+1)k\le2s+1<(2t+3)k$, hence $0\le b<k^2$ and $a>0$.

If $Y_{00}\ge0$, negate the identity and divide by
$k^2\norm{w'}^2$. The result is at most the convex combination, with
weights $a/k^2$ and $b/k^2$, of the component normalized violations.
At least one component is therefore no smaller. Repeating with the
greatest common divisor of the direction coefficients yields a
primitive direction. The condition on $Y_{00}$ is needed for this
domination conclusion; the algebraic identity itself has no such
restriction.
\end{proof}

\subsection{Rank-one examples for general integer moments}

The next results concern general integer moments. They do not describe
fractional rank-one Boolean moments: if $Y$ has rank one and also
$Y_{ii}=Y_{0i}$, then $x_i^2=x_i$, so $x\in\{0,1\}^n$ and every split is
satisfied.

\begin{proposition}\label{prop:rank-one-split-value}
Every rational symmetric rank-one $Y$ with $Y_{00}=1$ equals
$(1,x)(1,x)\trans$. Let $D$ be the least common denominator of the
coordinates of $x$, and put $p=D(1,x)\in\Z^{n+1}$. Then $\gcd(p)=1$ and
\[
 q_Y(v)=\frac{m(m+D)}{D^2},\qquad m=p\trans v.
\]
The point violates a split if and only if $D\ge2$. Its maximum raw
violation is $\lfloor D^2/4\rfloor/D^2$, and a maximizing split can be
found in polynomial time.
\end{proposition}

\begin{proof}
Rank one and $Y_{00}=1$ imply $Y=Y_{\cdot0}Y_{0\cdot}$. For any prime
dividing $D$, some reduced coordinate denominator contains its full
power in $D$, so the corresponding coordinate of $p$ is not divisible
by that prime. Hence $p$ is primitive, and $p\trans v$ ranges over all
integers. The displayed formula follows by expansion. It is negative
exactly when $-D<m<0$, and its minimum is attained at
$m=-\lfloor D/2\rfloor$ or $m=-\lceil D/2\rceil$. The extended Euclidean
algorithm finds a polynomial-size integer vector $a$ with
$p\trans a=1$; multiplying $a$ by the desired $m$ gives a maximizing
split. The least common denominator has polynomial encoding length.
\end{proof}

At fractional rational rank-one points this raw maximum is at least
$2/9$, even if $x$ is arbitrarily close to an integer vector. For example,
with one variable $x_1=1-1/D$, the elementary direction has raw violation
$(D-1)/D^2$, whereas $v=(-a,ae_1)$, $a=\lfloor D/2\rfloor$, attains the
raw maximum. For $D\ge4$ its normalized violation is
$(D-a)/(D^2a)<(D-1)/D^2$, so it is strictly smaller than that of the
elementary direction. This gives a reason to distinguish
raw violation from normalized violation; it is not a statement about
which choice gives the best objective bound.

\begin{proposition}\label{prop:restricted-rank-one-hardness}
At rational rank-one general integer moment points, deciding whether
there exists a violated split vector $v\in\{0,1\}^{n+2}$ is
$\NP$-complete, although
unrestricted split separation is polynomial there.
\end{proposition}

\begin{proof}
Given a SUBSET SUM instance with positive integers $a_1,\ldots,a_n,T$
\citep{Karp1972},
set $x_i=-12a_i/5$ and $x_{n+1}=(12T-2)/5$, and take
$Y=(1,x)(1,x)\trans$. For a $0/1$ vector $v$ put
$A_I=\sum_{i:v_i=1}a_i$. The split is violated if and only if
$-1<\tau<0$, where
$\tau=v_0+\sum_i v_ix_i$.
If $v_{n+1}=0$, then $\tau=v_0-12A_I/5$: it is nonnegative when
$A_I=0$ and at most $-7/5$ otherwise. If $v_{n+1}=1$, then
\[
 \tau=v_0-\frac25+\frac{12}{5}(T-A_I).
\]
For $T=A_I$, the choice $v_0=0$ gives $\tau=-2/5$. For $T-A_I\le-1$,
both offsets give $\tau\le-9/5$; for $T-A_I\ge1$, both give
$\tau\ge2$. Thus a restricted violated split exists exactly when the
subset-sum instance is solvable. The rational matrix has polynomial
encoding length, and restricted membership in $\NP$ is immediate.
This reduction does not claim strong hardness or hardness at Boolean
rank-one points.
\end{proof}
